\documentclass[UTF8,fontset=fandol,10pt]{ctexart}
\usepackage[a4paper,margin=19mm]{geometry}
\usepackage{amsmath,amssymb,array,longtable,booktabs,url,xcolor,fancyhdr}
\usepackage[colorlinks=true,linkcolor=blue,urlcolor=blue]{hyperref}
\pagestyle{fancy}\fancyhf{}\fancyhead[L]{YunTAPR-Net / synthetic boundaries}
\fancyhead[R]{2026-10-10 / v1}\fancyfoot[C]{\thepage}\setlength{\headheight}{14pt}
\setlength{\parindent}{0pt}\setlength{\parskip}{4pt}\setlength{\emergencystretch}{4em}
\begin{document}
\begin{center}\Large YunTAPR-Net：Phase-B v2\\合成边界验证与最终代码审查交接
\end{center}
\noindent\textbf{SYNTHETIC\_ENGINEERING\_ONLY；研究者待审，正式执行未授权。}
\section{正式集成前代码审查交接}

基线 ccfa6869672c\allowbreak{}96634fede157\allowbreak{}b06984741cc7\allowbreak{}1c1b。本轮完成剩余合成批次边界：49 CPU、12 CUDA、16静态项通过；原91 CPU/\allowbreak{}6 CUDA仅核对身份，未重跑。全部为 SYNTHETIC\_\allowbreak{}ENGINEERING\_\allowbreak{}ONLY，不能证明概率校准、q32欠覆盖或上尾问题已改善。审批状态没有升级。

\subsection{本人应审的准确源码}

路径前缀：A=src/\allowbreak{}yuntapr/\allowbreak{}experimental/\allowbreak{}phase\_\allowbreak{}b\_\allowbreak{}v2\_\allowbreak{}ablations/\allowbreak{}；I=src/\allowbreak{}yuntapr/\allowbreak{}experimental/\allowbreak{}phase\_\allowbreak{}b\_\allowbreak{}v2\_\allowbreak{}integration/\allowbreak{}；B=src/\allowbreak{}yuntapr/\allowbreak{}experimental/\allowbreak{}phase\_\allowbreak{}b\_\allowbreak{}v2\_\allowbreak{}boundaries/\allowbreak{}；M=src/\allowbreak{}yuntapr/\allowbreak{}models/\allowbreak{}quantile\_\allowbreak{}v2/\allowbreak{}；T=src/\allowbreak{}yuntapr/\allowbreak{}training/\allowbreak{}；F=src/\allowbreak{}yuntapr/\allowbreak{}losses/\allowbreak{}。每个函数起止行与完整SHA见 HANDOFF\_\allowbreak{}SOURCE\_\allowbreak{}MAP.json；下表是同一索引中的实际符号，全部需研究者阅读，不含审批勾选。

\begingroup\small\setlength{\tabcolsep}{3pt}\renewcommand{\arraystretch}{1.12}
\begin{longtable}{>{\raggedright\arraybackslash}p{\dimexpr(\linewidth-6\tabcolsep)/3\relax}>{\raggedright\arraybackslash}p{\dimexpr(\linewidth-6\tabcolsep)/3\relax}>{\raggedright\arraybackslash}p{\dimexpr(\linewidth-6\tabcolsep)/3\relax}}
\toprule
源文件 & 函数或类 & 本人审查重点 \\
\midrule\endfirsthead
源文件 & 函数或类 & 本人审查重点 \\
\midrule\endhead
A/\allowbreak{}config.py & get\_\allowbreak{}config, RunSpec, workload, learning\_\allowbreak{}rate\_\allowbreak{}prefix & 闭合E0/\allowbreak{}E1/\allowbreak{}E2；候选种子与18身份；50epoch轨迹前9epoch \\
A/\allowbreak{}validation.py & supervision, graph\_\allowbreak{}zero & FP32比较后提升；独立mask；无雨零图与零监督停止 \\
A/\allowbreak{}focal.py & occurrence\_\allowbreak{}numerator & 稳定BCE；gamma0仍保留alpha=.5 \\
A/\allowbreak{}pinball.py & frozen\_\allowbreak{}taus, conditional\_\allowbreak{}pinball\_\allowbreak{}numerator & 32tau均值、FP64、仅有雨有效像元、无物理转换 \\
A/\allowbreak{}total.py & combine\_\allowbreak{}numerators, candidate\_\allowbreak{}loss, CandidateLossResult & lambda仅乘一次；训练N\_\allowbreak{}valid和报告N\_\allowbreak{}rain \\
M/\allowbreak{}models.py & B0MatchedV2, B1V2 & 1/\allowbreak{}6帧完整501网格；原SP04；不能替换结构 \\
M/\allowbreak{}heads.py & ProbabilityHeadsV2 & BF16发生；FP32 raw33；FP64输出q32 \\
M/\allowbreak{}parameterization.py & normalized\_\allowbreak{}monotonic\_\allowbreak{}quantiles & 归一化正增量与独立span，epsilon不改 \\
M/\allowbreak{}outputs.py & validate\_\allowbreak{}log\_\allowbreak{}quantiles & 严格单调、支撑和非有限值停止 \\
I/\allowbreak{}initialization.py & seeded\_\allowbreak{}environment, fresh\_\allowbreak{}paired\_\allowbreak{}models & 独立reseed；同形复制；B1两输入权重保留；状态/\allowbreak{}非持久buffer身份 \\
I/\allowbreak{}controls.py & synthetic\_\allowbreak{}epoch\_\allowbreak{}order, endpoint\_\allowbreak{}boundary & seed+epoch独立排列；9epoch终点；无BEST替代 \\
I/\allowbreak{}controls.py & reject\_\allowbreak{}resume, reject\_\allowbreak{}formal\_\allowbreak{}start & 缺批准永拒；字符串不是批准；不加载LAST \\
A/\allowbreak{}readiness.py & review\_\allowbreak{}checkpoint\_\allowbreak{}binding, BlockedRunner & 具体run/\allowbreak{}LAST/\allowbreak{}code/\allowbreak{}data/\allowbreak{}init绑定；无启动能力 \\
I/\allowbreak{}resources.py /\allowbreak{} safety.py & resource\_\allowbreak{}snapshot, require\_\allowbreak{}resources, install\_\allowbreak{}process\_\allowbreak{}file\_\allowbreak{}guard, synthetic\_\allowbreak{}operation\_\allowbreak{}guards & 模型构造前准入；不足停止；进程守卫不是OS沙箱 \\
B/\allowbreak{}boundary.py & BoundaryBatch, make\_\allowbreak{}boundary\_\allowbreak{}pair, losses\_\allowbreak{}from\_\allowbreak{}output, forward\_\allowbreak{}boundary & 只允许train1、validation8/\allowbreak{}5；M1/\allowbreak{}Q1；inference\_\allowbreak{}mode \\
B/\allowbreak{}boundary.py & evaluation\_\allowbreak{}components, aggregate\_\allowbreak{}evaluation & 复用冻结FP64验证器；未加权先累加再相除 \\
B/\allowbreak{}gradients.py & checked\_\allowbreak{}gradients & 全部参数梯度；干批分位数头零梯度；NaN/\allowbreak{}缺失停止 \\
F/\allowbreak{}total\_\allowbreak{}loss.py & b0\_\allowbreak{}core\_\allowbreak{}loss & E0数学参照；零监督返回skip标记与正式入口停止对应 \\
T/\allowbreak{}phase\_\allowbreak{}a\_\allowbreak{}validation\_\allowbreak{}v2.py & LogDomainValidation & 固定gamma2，发生项FP64重算；不是训练S\_\allowbreak{}occ \\
T/\allowbreak{}formal\_\allowbreak{}phase\_\allowbreak{}a\_\allowbreak{}v2.py & paired\_\allowbreak{}initialization, forward\_\allowbreak{}loss, update, loader & 只读对照；训练batch2尾1，验证8尾5；update本轮未调用 \\
T/\allowbreak{}checkpoint\_\allowbreak{}v2.py & validate\_\allowbreak{}payload, verify\_\allowbreak{}file, apply\_\allowbreak{}verified, CheckpointStore & Phase-A 50epoch/\allowbreak{}BEST schema不能直接作B9/\allowbreak{}V0入口；后3符号未调用 \\
\bottomrule\end{longtable}\endgroup

\subsection{必须理解的公式与实际边界}

valid是两个掩膜交集；rain先在原FP32参考上严格比较，不先cast到FP64。人工3430格点仅测试计数，不代表真实云南地理配准。

\[
v=m_{\mathrm{reference}}\land m_{\mathrm{artificial}},\quad
z=\mathbf1\{y_{32}>\operatorname{float32}(0.1)\},\quad
N_v=\sum v,\quad N_r=\sum vz.
\]


Focal映射A/\allowbreak{}focal.py。E0/\allowbreak{}E2的gamma=2，E1的gamma=0；所有alpha=.5。BF16 logit的BCE提升FP32，gamma0不等于未加权BCE。

\[
b=\operatorname{BCEWithLogits}(\ell,z),\quad p_t=e^{-b},\quad
S_{\rm occ}=\sum_v\alpha_t(1-p_t)^\gamma b,\quad
\alpha_t=\alpha z+(1-\alpha)(1-z).
\]


q映射M/\allowbreak{}parameterization.py：32allocation+1span原通道，正权重和FP64前缀；不是排序或截断输出。两个epsilon均1e-4，qlog支撑/\allowbreak{}单调检查仍由原函数执行。

\[
w_i=\operatorname{softplus}(a_i)+\epsilon_w,\quad
q_i=\log(1+0.1)+(\operatorname{softplus}(s)+\epsilon_s)
\frac{\sum_{j\le i}w_j}{\sum_{j=1}^{32}w_j},\quad
\tau_i=\frac{i-0.5}{32}.
\]


Pinball映射A/\allowbreak{}pinball.py与A/\allowbreak{}total.py；先均值32tau，再仅在有雨有效像元求和。lambda不进入科学Conditional Pinball。空雨时S\_\allowbreak{}qr=0、图连接、CPB=None；零有效监督停止。

\[
S_{\rm qr}=\sum_{vz}\frac1{32}\sum_{i=1}^{32}
\max\{\tau_i e_i,(\tau_i-1)e_i\},\quad
e_i=\log(1+y_{32})-q_i,
\]

\[
L_{\rm train}=\frac{S_{\rm occ}+\lambda_qS_{\rm qr}}{N_v},\qquad
\mathrm{CPB}=\frac{\sum S_{\rm qr}}{\sum N_r}\ \ (\sum N_r>0).
\]


共同Core必须从T/\allowbreak{}phase\_\allowbreak{}a\_\allowbreak{}validation\_\allowbreak{}v2.py在FP64重算固定gamma2发生分子并全局累加，使用lambda=1的未加权分位数项。不能将E1/\allowbreak{}E2的training\_\allowbreak{}objective当共同Core，甚至E0训练发生分子也不是同精度的验证分子。本轮首次静态发现并修复这一差别，旧attempt完整保留。

\subsection{尚未通过或未执行的检查}

本轮定义的49 CPU/\allowbreak{}12 CUDA边界检查均通过，没有待修的数值失败；下列是未执行项，不能写成通过。

\noindent\textbullet\quad 真实2023/\allowbreak{}2024数据资格、冻结地理mask/\allowbreak{}SP04配准、shared scaler实际应用、真实ID顺序和读取防火墙。需要独立真实数据preflight许可，2025继续封存。\par
\noindent\textbullet\quad 正式optimizer/\allowbreak{}clip/\allowbreak{}scheduler更新事务、长时epoch循环、实际吞吐与共享GPU持续资源风险。没有step；无优化器显存测量。\par
\noindent\textbullet\quad 新v2 runner、B9终点checkpoint schema、原子写入/\allowbreak{}耐久性、故障注入和审批验真。没有私有checkpoint加载、保存或恢复。\par
\noindent\textbullet\quad 新批次边界本轮仅seed2026；另外两seed初始化与训练顺序核验使用原始已通过证据，不伪称重新覆盖。CPU完整模型forward/\allowbreak{}backward未运行，实际完整路径在CUDA执行。\par
\noindent\textbullet\quad 正式未来runner中异常停止、阶段批准、LAST绑定审批和失败后人工处置的端到端测试；当前只有隔离保护与静态候选，不能启动训练。\par

\subsection{必须本人决定的科学与治理问题}

全部保持 RESEARCHER\_\allowbreak{}DECISION\_\allowbreak{}REQUIRED：H-O/\allowbreak{}H-Q及E0/\allowbreak{}E1/\allowbreak{}E2最终科学接受；三个seed、D1/\allowbreak{}Q1/\allowbreak{}N0/\allowbreak{}I3/\allowbreak{}B9/\allowbreak{}S0/\allowbreak{}V0；Brier与q32覆盖误差主指标、完整32tau与强雨Pinball/\allowbreak{}上尾副作用；最小有意义效应、副作用界限、多重比较仍NOT\_\allowbreak{}YET\_\allowbreak{}ESTABLISHED；2024历史开发验证的解释范围，不能作为全新独立确认集。

Phase-A科学接受范围和B1历史恢复处置须独立决定，历史恢复仍NOT\_\allowbreak{}GRANTED。数据权限、计算资源、正式集成范围、具体阶段执行及LAST恢复均另需明确授权。文档、JSON、测试PASS和Git commit不构成批准事件。

建议本人按 A/\allowbreak{}config→validation→focal/\allowbreak{}pinball/\allowbreak{}total→M/\allowbreak{}heads/\allowbreak{}parameterization→I/\allowbreak{}initialization/\allowbreak{}controls→B/\allowbreak{}boundary/\allowbreak{}gradients→测试顺序阅读；现场解释FP32阈值反例、两个分母、空雨零梯度、FP64共同验证、epoch9终点与LAST绑定。随后审查正式集成差异与独立执行授权。本轮在审批边界停止，不派生新工作包。

V2\_\allowbreak{}PHASE\_\allowbreak{}B\_\allowbreak{}AUTHORIZED=false；RESEARCHER\_\allowbreak{}SCIENTIFIC\_\allowbreak{}APPROVAL\_\allowbreak{}REQUIRED=true；HISTORICAL\_\allowbreak{}RECOVERY\_\allowbreak{}RATIFICATION=NOT\_\allowbreak{}GRANTED；2025\_\allowbreak{}RAW\_\allowbreak{}ACCESS=2025\_\allowbreak{}PIXELS\_\allowbreak{}READ=0。
\clearpage\section{合成边界测试与资源实测}

全部数据为人工有限张量；人工mask每场3430格点、M1全部原生像元有效、Q1参考完整。场景ID显式SYNTHETIC\_\allowbreak{}ENGINEERING\_\allowbreak{}ONLY，没有真实年份或路径。B0末帧取自B1六帧；模型、SP04公开映射、tau与epsilon保持原样。

最终CPU49、CUDA12、静态16通过，0失败/\allowbreak{}错误/\allowbreak{}跳过/\allowbreak{}warning。唯一pytest覆盖61；初次与复测共122个实例，不能相加作独立证据。旧91/\allowbreak{}6未重跑。两次各12完整forward/\allowbreak{}12backward，累计24/\allowbreak{}24；最终4个train1案例各3臂backward，8个validation8/\allowbreak{}5案例仅inference\_\allowbreak{}mode。

\begingroup\small\setlength{\tabcolsep}{3pt}\renewcommand{\arraystretch}{1.12}
\begin{longtable}{>{\raggedright\arraybackslash}p{\dimexpr(\linewidth-12\tabcolsep)/6\relax}>{\raggedright\arraybackslash}p{\dimexpr(\linewidth-12\tabcolsep)/6\relax}>{\raggedright\arraybackslash}p{\dimexpr(\linewidth-12\tabcolsep)/6\relax}>{\raggedright\arraybackslash}p{\dimexpr(\linewidth-12\tabcolsep)/6\relax}>{\raggedright\arraybackslash}p{\dimexpr(\linewidth-12\tabcolsep)/6\relax}>{\raggedright\arraybackslash}p{\dimexpr(\linewidth-12\tabcolsep)/6\relax}}
\toprule
batch & 语义 & N\_\allowbreak{}valid & allocated MiB & reserved MiB & forward+三臂loss ms \\
\midrule\endfirsthead
batch & 语义 & N\_\allowbreak{}valid & allocated MiB & reserved MiB & forward+三臂loss ms \\
\midrule\endhead
1 & train forward/\allowbreak{}backward & 3430 & 1435.02–1451.44 & 1820.00–1832.00 & 55.45–432.72 \\
8 & validation inference & 27440 & 3280.71–3329.14 & 4254.00–4628.00 & 309.63–317.92 \\
5 & validation inference & 17150 & 2057.67–2087.23 & 2900.00–2916.00 & 197.74–211.05 \\
\bottomrule\end{longtable}\endgroup

GPU为RTX5060 Laptop、PyTorch2.11.0+cu128、CUDA12.8。每例先要求主机可用\(\ge\)8GiB、CUDA API free\(\ge\)5GiB、工作盘\(\ge\)1GiB、BF16可用，再构造原模型；验证也保守复用此门槛，snapshot的full\_\allowbreak{}backward=true表示资源门槛选择，不表示验证做了backward。没有OOM、超时、降batch/\allowbreak{}精度/\allowbreak{}结构、换seed或自动重试。

每值仅一批，首例冷启动；cuda synchronize+perf\_\allowbreak{}counter计时。allocator峰值包括已有合成输入、三臂loss/\allowbreak{}梯度及共同FP64评价检查，不含其他进程、AdamW状态、真实I/\allowbreak{}O、checkpoint事务或長时循环。显存准入不是容量预测，旁路进程仍会影响可用资源。

E0数值与原b0\_\allowbreak{}core\_\allowbreak{}loss在冻结BF16环境比较，12例通过既有rtol=8e-7/\allowbreak{}atol=1e-8；最大绝对差为 1.8359273240564633e-08。本轮没有重复旧batch2全参数梯度差异检查；新batch1三臂全部74参数梯度存在、形状/\allowbreak{}dtype/\allowbreak{}有限性正确。混合有雨两头有信号；空雨分位数头梯度精确为0且不缺失，科学CPB=None。无参数改变，各例初始SHA与原seed2026证据一致。

训练N\_\allowbreak{}valid分别3430/\allowbreak{}27440/\allowbreak{}17150；混合雨N\_\allowbreak{}rain分别2450/\allowbreak{}19600/\allowbreak{}12250，仅为人工计数。CPU检查掩膜外NaN参考只在低层loss合成夹具中排除，不能视为M1/\allowbreak{}Q1合格模型场景。完整模型遇到缺原生帧、缺监督、非有限或身份/\allowbreak{}精度异常必须停止。

所有分子/\allowbreak{}目标数值保存在 SYNTHETIC\_\allowbreak{}BOUNDARY\_\allowbreak{}RESULTS.json；这是工程算术而非新增降水性能指标。原始case/\allowbreak{}资源/\allowbreak{}源码快照见tests/\allowbreak{}cuda\_\allowbreak{}attempt\_\allowbreak{}002\_\allowbreak{}evidence/\allowbreak{}。共同科学评价按冻结FP64验证器，训练objective仍各臂独立，不用于共同Core。

初次CPU49/\allowbreak{}CUDA12通过后，静态发现共同评价接口误用E0 FP32训练发生分子；已改为复用冻结FP64 LogDomainValidation，增加精确source对照，复测49/\allowbreak{}12/\allowbreak{}16通过，容差未放宽。初次测试、source index、static16之前的static15与修复说明均保留，不伪造数值失败。详见 修复记录。
\clearpage\section{新增模块审查笔记}

本轮仅3个候选Python文件与3个新测试文件；不重写旧损失或旧适配器。完整旧损失学习见 上一轮学习材料。

\subsection{boundary.py}

verify\_\allowbreak{}inherited\_\allowbreak{}identity只核对226个既有公开文件字节；SHA说明身份，不能证明独立科学批准。BoundaryBatch是frozen dataclass，字段名固定但Tensor内容仍可变，所以每次validate。x为FP32[B,1或6,501,501]；参考FP32[B,1,100,100]；三个bool mask独立且同设备。只接受B=1/\allowbreak{}8/\allowbreak{}5和相应phase，data\_\allowbreak{}years必须为空。

make\_\allowbreak{}boundary\_\allowbreak{}pair用arange/\allowbreak{}reshape/\allowbreak{}sin生成有限值，expand后clone使人工mask可独立检查。nextafter生成float32(0.1)上下相邻值；这是阈值教学，不是读真实降雨。B0使用最后一帧clone，明确避免输入别名。to用dataclasses.replace返回转设备副本，不能把frozen解释为深度不可变。

losses\_\allowbreak{}from\_\allowbreak{}output校验BF16发生与FP64 qlog、sigmoid一致、shape/\allowbreak{}有限性；沿用get\_\allowbreak{}config/\allowbreak{}candidate\_\allowbreak{}loss，返回三臂而非搜索参数。forward\_\allowbreak{}boundary要求原模型精确类型和train/\allowbreak{}eval状态；with inference\_\allowbreak{}mode(mode=not training)控制计算图；with autocast只控制冻结BF16路径，内部raw33与qlog升精度仍由原头负责。函数不做backward或step。

evaluation\_\allowbreak{}components使用原LogDomainValidation重新算FP64发生分子；不是detach训练FP32分子后冒称FP64。aggregate\_\allowbreak{}evaluation先fsum未加权分子与分母，再求比值；不能平均batch均值或把空雨CPB填0。lambda与gamma变化只影响训练报告。

\subsection{gradients.py与包入口}

checked\_\allowbreak{}gradients逐一检查parameter.grad；backward累加梯度，所以测试每臂前zero\_\allowbreak{}grad(set\_\allowbreak{}to\_\allowbreak{}none=True)。这只是清梯度，不更新参数。train1同一输出图供三臂使用，前两次retain\_\allowbreak{}graph=True，最后释放；这样只改变loss，不能声称模拟三个正式训练run。空雨graph\_\allowbreak{}zero沿计算图传播0，允许梯度全0，但不允许grad缺失或NaN。包入口仅scope常量，没有runner。

\subsection{测试与工具}

test\_\allowbreak{}cpu\_\allowbreak{}boundaries用完整目标网格核对两个分母、float32阈值、掩膜、异常与不等批次汇总；不是CPU完整模型训练。test\_\allowbreak{}cuda\_\allowbreak{}boundaries每例先资源准入、再fresh原模型，再forward，train1才backward；原始权重、全参数梯度和buffer身份只存SHA/\allowbreak{}范数，不保存模型。conftest复用旧step/\allowbreak{}load/\allowbreak{}save/\allowbreak{}expm1保护；Python守卫可被进程外代码改写，不能成为正式授权认证。

run\_\allowbreak{}checks调用新测试子进程，保存原始私有log及脱敏公开XML、执行时源码SHA，不自动重试。static\_\allowbreak{}review只AST/\allowbreak{}文件哈希/\allowbreak{}已发表记录，不import模型。build\_\allowbreak{}review/\allowbreak{}export/\allowbreak{}close只处理文档与身份；doc自身成功不构成批准。

研究者小练习：解释为何float32(0.1)先double再比较Python0.1可能错标有雨；手算batch1的3430分母；解释空雨q头zero与缺失梯度的区别；用8+5批次证明先累加分子、再相除；找出Phase-A checkpoint的50epoch/\allowbreak{}BEST字段为何不能直接替代B9/\allowbreak{}V0。
\end{document}
